\documentclass[12pt,oneside,a4paper]{article}
\pagenumbering{arabic}
\usepackage{graphicx}
\title{Quiz primo compitino}
\date{17/04/2024}
\begin{document}
\maketitle
\textbf{Ho inserito in questo file LaTeX prima gli esercizi dei quiz in modo che vi possiate esercitare, le soluzioni sono inserite in fondo al pdf}
%h serve per indicare che l'immagine vada "here", t per "top", b per "bottom", p per "on an extra page" e ! per "ovverride" che forza la locazione specifica
\newpage
\section{Esercizi}
\subsection{Esercizio 1}
\begin{figure}[h]
\includegraphics[scale=0.9]{Esuno.png}
\end{figure}
\newpage
\subsection{Esercizio 2}
\begin{figure}[h] 
\includegraphics[scale=0.9]{Esdue.png}
\end{figure}
\newpage
\subsection{Esercizio 3}
\begin{figure}[h]
\includegraphics[scale=0.85]{Estre.png}
\end{figure}
\newpage
\subsection{Esercizio 4}
\begin{figure}[h]
\includegraphics[scale=0.9]{Esquattro.png}
\end{figure}
\newpage
\section{Soluzioni}
\subsection{Soluzione esercizio 1}
PRE: n$>$=1\newline
POST: restituisce la potenza di 2 elevato alla n-1
\subsection{Soluzione esercizio 2}
//1: stampa 0 in quanto il puntatore p punta ad x=2, modificando il suo valore nella riga precedente per via dell'aritmetica del C, la quale esegue prima 2/4=0 (troncando il valore ottenuto) moltiplicandolo per 2 ottenendo 0 nuovamente\newline
//2: dangling reference, l'indice dell'array è indicato dal puntatore p il quale però punta ad una variabile deallocata, quindi indefinita\newline
//3: stampa 6: infatti *(m+2) punta al primo elemeno della seconda riga ma con (*(m+2)[1]) ci si riferisce al secondo elemento della seconda riga, ovvero 6
\subsection{Soluzione esercizio 3}
char *x[5]\newline
double* (*p)[7]
\subsection{Soluzione esercizio 4}
INV1
\end{document}